\documentclass{article}

\def\npart {IA}
\def\nterm {Lent}
\def\nyear {2024}
\def\nlecturer {Prof. J.\ Norris}
\def\ncourse {Probability}

\setcounter{tocdepth}{2}

\usepackage{alltt}
\usepackage{amsmath, amsfonts, amssymb, amsthm}
\usepackage{booktabs}
\usepackage{caption}
\usepackage{enumitem}
\usepackage{fancyhdr}
\usepackage{graphicx}
\usepackage{mathdots}
\usepackage{mathtools}
\usepackage{microtype}
\usepackage{multirow}
\usepackage{pdflscape}
\usepackage{pgfplots}
\usepackage{siunitx}
\usepackage{textcomp}
\usepackage{slashed}
\usepackage{tabularx}
\usepackage{tikz}
\usepackage{tkz-euclide}
\usepackage[normalem]{ulem}
\usepackage[all]{xy}

\pgfplotsset{compat=1.18}
\pagestyle{fancyplain}

\lhead{\emph{\nouppercase{\leftmark}}}
\rhead{
  \ifnum\thepage=1
  \else
    \npart\ \vline\ \ncourse
  \fi}

\def\nauthor{Marcus Ng}
\author{Based on lectures by \nlecturer \\\small Notes taken by \nauthor}
\date{\nterm\ \nyear}
\title{Part \npart\ --- \ncourse}

\newcommand*{\Cdot}{{\raisebox{-0.25ex}{\scalebox{1.5}{$\cdot$}}}}
\setlist[enumerate,1]{label={(\roman*)}}

\newcommand {\pd}[2][ ]{
  \ifx #1 { }
    \frac{\partial}{\partial #2}
  \else
    \frac{\partial^{#1}}{\partial #2^{#1}}
  \fi
}

% Theorems
\theoremstyle{definition}
\newtheorem*{aim}{Aim}
\newtheorem*{axiom}{Axiom}
\newtheorem*{claim}{Claim}
\newtheorem*{cor}{Corollary}
\newtheorem*{conjecture}{Conjecture}
\newtheorem*{defi}{Definition}
\newtheorem*{eg}{Example}
\newtheorem*{ex}{Exercise}
\newtheorem*{fact}{Fact}
\newtheorem*{law}{Law}
\newtheorem*{lemma}{Lemma}
\newtheorem*{notation}{Notation}
\newtheorem*{prop}{Proposition}
\newtheorem*{question}{Question}
\newtheorem*{problem}{Problem}
\newtheorem*{rrule}{Rule}
\newtheorem*{thm}{Theorem}
\newtheorem*{assumption}{Assumption}
\newtheorem*{assert}{Assertion}

\newtheorem*{remark}{Remark}
\newtheorem*{warning}{Warning}
\newtheorem*{exercise}{Exercise}

\newtheorem{nthm}{Theorem}[section]
\newtheorem{nlemma}[nthm]{Lemma}
\newtheorem{nprop}[nthm]{Proposition}
\newtheorem{ncor}[nthm]{Corollary}

\SetEnumitemKey{cases}{
  label=\underline{Case~\arabic*:},
  ref={\arabic*},
  align=left,
  leftmargin=\parindent
}

% \renewcommand{\labelitemi}{--}
% \renewcommand{\labelitemii}{$\circ$}
% \renewcommand{\labelenumi}{(\roman{*})}

\let\stdsection\section
\renewcommand\section{\newpage\stdsection}

% Strike through
\def\st{\bgroup \ULdepth=-.55ex \ULset}


%%%%%%%%%%%%%%%%%%%%%%%%%
%%%%% Maths Symbols %%%%%
%%%%%%%%%%%%%%%%%%%%%%%%%

% Logical symbols
\newcommand{\contradiction}{
  \begin{tikzpicture}[rotate=45,x=0.5ex,y=0.5ex]
  \draw[line width=.1ex] (0,2) -- (3,2) (0,1) -- (3,1) (1,3) -- (1,0) (2,3) -- (2,0);
  \end{tikzpicture}
}


\let\U\relax
\let\C\relax
\let\G\relax

% Matrix groups
% \newcommand{\GL}{\mathrm{GL}}
% \newcommand{\Or}{\mathrm{O}}
% \newcommand{\PGL}{\mathrm{PGL}}
% \newcommand{\PSL}{\mathrm{PSL}}
% \newcommand{\PSO}{\mathrm{PSO}}
% \newcommand{\PSU}{\mathrm{PSU}}
% \newcommand{\SL}{\mathrm{SL}}
% \newcommand{\SO}{\mathrm{SO}}
% \newcommand{\Spin}{\mathrm{Spin}}
% \newcommand{\Sp}{\mathrm{Sp}}
% \newcommand{\SU}{\mathrm{SU}}
\newcommand{\U}{\mathrm{U}}
% \newcommand{\Mat}{\mathrm{Mat}}

% Matrix algebras
% \newcommand{\gl}{\mathfrak{gl}}
% \newcommand{\ort}{\mathfrak{o}}
% \newcommand{\so}{\mathfrak{so}}
% \newcommand{\su}{\mathfrak{su}}
% \newcommand{\uu}{\mathfrak{u}}
% \renewcommand{\sl}{\mathfrak{sl}}

% Special sets
\newcommand{\C}{\mathbb{C}}
\newcommand{\CP}{\mathbb{CP}}
\newcommand{\GG}{\mathbb{G}}
\newcommand{\N}{\mathbb{N}}
\newcommand{\Q}{\mathbb{Q}}
\newcommand{\R}{\mathbb{R}}
\newcommand{\RP}{\mathbb{RP}}
\newcommand{\T}{\mathbb{T}}
\newcommand{\Z}{\mathbb{Z}}
\renewcommand{\H}{\mathbb{H}}

% Brackets
\newcommand{\abs}[1]{\left\lvert #1\right\rvert}
\newcommand{\bket}[1]{\left\lvert #1\right\rangle}
\newcommand{\brak}[1]{\left\langle #1 \right\rvert}
\newcommand{\braket}[2]{\left\langle #1\middle\vert #2 \right\rangle}
\newcommand{\bra}{\langle}
\newcommand{\ket}{\rangle}
\newcommand{\norm}[1]{\left\lVert #1\right\rVert}
\newcommand{\normalorder}[1]{\mathop{:}\nolimits\!#1\!\mathop{:}\nolimits}
\newcommand{\tv}[1]{|#1|}
\renewcommand{\vec}[1]{\boldsymbol{\mathbf{#1}}}

% not-math
% \newcommand{\bolds}[1]{{\bfseries #1}}
% \newcommand{\cat}[1]{\mathsf{#1}}
% \newcommand{\ph}{\,\cdot\,}
% \newcommand{\term}[1]{\emph{#1}\index{#1}}
% \newcommand{\phantomeq}{\hphantom{{}={}}}

% Probability
% \DeclareMathOperator{\Bernoulli}{Bernoulli}
% \DeclareMathOperator{\betaD}{beta}
% \DeclareMathOperator{\bias}{bias}
% \DeclareMathOperator{\binomial}{binomial}
% \DeclareMathOperator{\corr}{corr}
% \DeclareMathOperator{\cov}{cov}
% \DeclareMathOperator{\gammaD}{gamma}
% \DeclareMathOperator{\mse}{mse}
% \DeclareMathOperator{\multinomial}{multinomial}
% \DeclareMathOperator{\Poisson}{Poisson}
% \DeclareMathOperator{\var}{var}
% \newcommand{\E}{\mathbb{E}}
% \newcommand{\Prob}{\mathbb{P}}

% Algebra
% \DeclareMathOperator{\adj}{adj}
% \DeclareMathOperator{\Ann}{Ann}
% \DeclareMathOperator{\Aut}{Aut}
% \DeclareMathOperator{\Char}{char}
% \DeclareMathOperator{\disc}{disc}
% \DeclareMathOperator{\dom}{dom}
% \DeclareMathOperator{\fix}{fix}
% \DeclareMathOperator{\Hom}{Hom}
% \DeclareMathOperator{\id}{id}
% \DeclareMathOperator{\image}{image}
% \DeclareMathOperator{\im}{im}
% \DeclareMathOperator{\re}{re}
% \DeclareMathOperator{\tr}{tr}
% \DeclareMathOperator{\Tr}{Tr}
% \newcommand{\Bilin}{\mathrm{Bilin}}
% \newcommand{\Frob}{\mathrm{Frob}}

% Others
% \newcommand\ad{\mathrm{ad}}
% \newcommand\Art{\mathrm{Art}}
% \newcommand{\B}{\mathcal{B}}
% \newcommand{\cU}{\mathcal{U}}
% \newcommand{\Der}{\mathrm{Der}}
% \newcommand{\D}{\mathrm{D}}
% \newcommand{\dR}{\mathrm{dR}}
% \newcommand{\exterior}{\mathchoice{{\textstyle\bigwedge}}{{\bigwedge}}{{\textstyle\wedge}}{{\scriptstyle\wedge}}}
% \newcommand{\F}{\mathbb{F}}
% \newcommand{\G}{\mathcal{G}}
% \newcommand{\Gr}{\mathrm{Gr}}
% \newcommand{\haut}{\mathrm{ht}}
% \newcommand{\Hol}{\mathrm{Hol}}
% \newcommand{\hol}{\mathfrak{hol}}
% \newcommand{\Id}{\mathrm{Id}}
% \newcommand{\lie}[1]{\mathfrak{#1}}
% \newcommand{\op}{\mathrm{op}}
% \newcommand{\Oc}{\mathcal{O}}
% \newcommand{\pr}{\mathrm{pr}}
% \newcommand{\Ps}{\mathcal{P}}
% \newcommand{\pt}{\mathrm{pt}}
% \newcommand{\qeq}{\mathrel{``{=}"}}
% \newcommand{\Rs}{\mathcal{R}}
% \newcommand{\Vect}{\mathrm{Vect}}
% \newcommand{\wsto}{\stackrel{\mathrm{w}^*}{\to}}
% \newcommand{\wt}{\mathrm{wt}}
% \newcommand{\wto}{\stackrel{\mathrm{w}}{\to}}
% \renewcommand{\d}{\mathrm{d}}
% \renewcommand{\P}{\mathbb{P}}
%\renewcommand{\F}{\mathcal{F}}

\let\Im\relax
\let\Re\relax

% \DeclareMathOperator{\area}{area}
% \DeclareMathOperator{\card}{card}
% \DeclareMathOperator{\ccl}{ccl}
% \DeclareMathOperator{\ch}{ch}
% \DeclareMathOperator{\cl}{cl}
% \DeclareMathOperator{\cls}{\overline{\mathrm{span}}}
% \DeclareMathOperator{\coker}{coker}
% \DeclareMathOperator{\conv}{conv}
% \DeclareMathOperator{\cosec}{cosec}
% \DeclareMathOperator{\cosech}{cosech}
% \DeclareMathOperator{\covol}{covol}
% \DeclareMathOperator{\diag}{diag}
% \DeclareMathOperator{\diam}{diam}
% \DeclareMathOperator{\Diff}{Diff}
% \DeclareMathOperator{\End}{End}
% \DeclareMathOperator{\energy}{energy}
% \DeclareMathOperator{\erfc}{erfc}
% \DeclareMathOperator{\erf}{erf}
% \DeclareMathOperator*{\esssup}{ess\,sup}
% \DeclareMathOperator{\ev}{ev}
% \DeclareMathOperator{\Ext}{Ext}
% \DeclareMathOperator{\fst}{fst}
% \DeclareMathOperator{\Fit}{Fit}
% \DeclareMathOperator{\Frac}{Frac}
% \DeclareMathOperator{\Gal}{Gal}
% \DeclareMathOperator{\gr}{gr}
% \DeclareMathOperator{\hcf}{hcf}
\DeclareMathOperator{\Im}{Im}
% \DeclareMathOperator{\Ind}{Ind}
% \DeclareMathOperator{\Int}{Int}
% \DeclareMathOperator{\Isom}{Isom}
% \DeclareMathOperator{\lcm}{lcm}
% \DeclareMathOperator{\length}{length}
% \DeclareMathOperator{\Lie}{Lie}
% \DeclareMathOperator{\like}{like}
% \DeclareMathOperator{\Lk}{Lk}
% \DeclareMathOperator{\Maps}{Maps}
% \DeclareMathOperator{\orb}{orb}
% \DeclareMathOperator{\ord}{ord}
% \DeclareMathOperator{\otp}{otp}
% \DeclareMathOperator{\poly}{poly}
% \DeclareMathOperator{\rank}{rank}
% \DeclareMathOperator{\rel}{rel}
% \DeclareMathOperator{\Rad}{Rad}
\DeclareMathOperator{\Re}{Re}
% \DeclareMathOperator*{\res}{res}
% \DeclareMathOperator{\Res}{Res}
% \DeclareMathOperator{\Ric}{Ric}
% \DeclareMathOperator{\rk}{rk}
% \DeclareMathOperator{\Rees}{Rees}
% \DeclareMathOperator{\Root}{Root}
% \DeclareMathOperator{\sech}{sech}
% \DeclareMathOperator{\sgn}{sgn}
% \DeclareMathOperator{\snd}{snd}
% \DeclareMathOperator{\Spec}{Spec}
% \DeclareMathOperator{\spn}{span}
% \DeclareMathOperator{\stab}{stab}
% \DeclareMathOperator{\St}{St}
% \DeclareMathOperator{\supp}{supp}
% \DeclareMathOperator{\Syl}{Syl}
% \DeclareMathOperator{\Sym}{Sym}
% \DeclareMathOperator{\vol}{vol}
\numberwithin{equation}{section}

\begin{document}
\maketitle{
    \small
    \noindent\textbf{Basic concepts}\\
    Classical probability, equally likely outcomes. Combinatorial analysis, permutations and combinations.
    Striling's formula (asymptotics for $\log n!$ proved).\hspace*{\fill}[3]
    
    \vspace{10pt}
    \noindent\textbf{Axiomatic approach}\\
    Axioms (countable case). Probability spaces. Inclusion-exclusion formula. Countinuity and subadditivity of probability measures. 
    Independance. Binomial, Poisson and geometric distributions. Relation between Poisson and binomial distributions. 
    Condition probability, Baye's formula. Examples, including Simpon's paradox.\hspace*{\fill}[5]

    \vspace{10pt}
    \noindent\textbf{Discrete random variables}\\
    Expectation. Functions of random variable, indicator function, variance, standard deviation.
    Covariance, independence of random variables. Generating functions: sums of independent random variables, random sum formula, moments.
    
    \vspace{5pt}
    \noindent Conditional expectation. Random walks: gambler's ruin, recurrence relations.
    Difference equations and their solution. Mean time to absorption, Branching processes: generating functions and extinction probability.
    Combinatorial applications of generating functions.\hspace*{\fill}[7]
    
    \vspace{10pt}
    \noindent\textbf{Continuous random variables}\\
    Distributions and density functions. Expectations: expectation of a function of a random variable.
    Uniform, normal and exponential random variables. Memoryless property of exponential distribution.

    \vspace{5pt}
    \noindent Joint distributions: transformation of random variables (including Jacobians), examples.
    SimulationL generating continuous random variables, Box-Muller transform, rejection sampling.
    Geometrical probability: Betrand's paradox, Buffon's needle. Correlation coefficient, bivariate normal random variables.\hspace*{\fill}[6]

    \vspace{10pt}
    \noindent\textbf{Inequalities and Limits}\\
    Markov's inequality, Chebyshev's inequality. Weak law of large numbers. 
    Convextivity: Jensen's inequality for general random variables, AM/GM inequality.

    \vspace{5pt}
    \noindent Moment generating functions and statement (no proof) of continuity theorem. 
    Statement of central limit theorem and sketch of proof. Examples, including sampling.\hspace*{\fill}[3]
}
\tableofcontents

%%%%%%%%%%%%%%%%%%%%%%
% Lecture 1: 19/01/2024
%%%%%%%%%%%%%%%%%%%%%%

\section{Mathematical Models for randomness}
\subsection{General definitions}

Let $\Omega$ be a set. Let $\F$ be a set of subsets of $\Omega$.
We say that $\F$ is a $\sigma$-algebra if
\begin{itemize}
    \item $\Omega \in \F$
    \item $\forall A \in \F$ and all sequences $(A_n : n \in \N)$ in $\F$
    \[\Rightarrow A^c \in \F, \qquad \bigcup_n A_n \in \F\]
\end{itemize}

Thus equivalently, $\F$ is non-empty and closed under countable set operations.
Note that 
\[ \bigcap_n A_n = \bkt{\bigcup_n A_n^c}^c\]
Let $\F$ be a $\sigma$-algebra. We say that a function
\[ \Prob: \F \rightarrow [0, 1] \]
is a \emph{probability measure} if
\begin{enumerate}
    \item $\Prob(\Omega) = 1$
    \item $\forall$ sequences $(A_n : n \in \N)$ of disjoint elements of $\F$,
    \[\Rightarrow \Prob \bkt{\bigcup_n A_n} = \sum_n \Prob(A_n)\]
    This is known as \emph{countable additivity}
\end{enumerate}

Let $\Prob$ be a probability measure on $(\Omega, \F)$, then $(\Omega, \F, \Prob)$ is a \emph{probability space}.

Elements of $\Omega$ are called \emph{outcomes} and elements of $\F$ are \emph{events}.
We interprit $\Prob(A)$ as the \emph{probability} of event $A$.

\subsection{Equally likely outcomes}
Let $\Omega$ be a finite set, we'll always take $\F$ to be $\wp(\Omega)$ in this case (unless stated).
Sometimes symmetry suggests a model with all outcomes equally likely,
\[\Prob(A) = \frac{\abs{A}}{\abs{\Omega}}\]

\begin{ex}
    Show that $(\Omega, \F, \Prob)$ is a probability space.
    This comes down to $\abs{A \cup B} = \abs{A} + \abs{B}$ for $A$, $B$ disjoint.
\end{ex}

\subsection{Throwing a dice}
$\Omega = {1, 2, 3, 4, 5, 6}$, $\F = \wp(\Omega)$ and $\Prob(A) = \frac{\abs{A}}{\abs{\Omega}}$.
For example,
\[\Prob({2, 5, 6}) = \frac{1}{2}\]

\subsection{Balls from a bag}
Let the balls be numbered $1, 2, \cdots, n$. Randomly remove all at once $k$ balls. 
What is the probability that we have chosen balls $1, 2, \cdots, k$.
We have 
\begin{enumerate}
    \item $\Omega = \{k \text{ element subsets of } \{1, \cdots, n\}\} \Rightarrow \abs{\Omega} = \binom{n}{k}$
    \item $A = \{\{1, \cdots, k\}\} \Rightarrow \abs{A} = 1$
\end{enumerate}
Hence,
\[\Prob(A) = \frac{1}{\abs{\Omega} = \binom{n}{k}^{-1}}\]

\subsection{Well-Shuffled pack of cards}
$\Omega =$ set of permutations of cards $\{1, 2, \cdots, 52\} \Rightarrow \abs{\Omega} = 52!$.
What is the probability for a well-shuffled pack that the first two cards are aces?
\begin{align*}
    \abs{A} &= 4 \times 3 \times 50! 
    \therefore \Prob(A) &= \frac{4 \times 3 \times 50! }{52!} = \frac{1}{221}
\end{align*}

\subsection{Distribution of the largest digit}
Observe from a string of random digits $n$ digits. What is the probability that the greatest of these digits is $k$?
\begin{align*}
    \Omega &= \{0, 1, \cdots, a\}^n  \Rightarrow \abs{\Omega} = 10^n\\
    A_k &= \{\text{all digits are } \leq k\}  \Rightarrow \abs{A_k} = (k+1)^n\\
    B_k &= \{\text{greatest digit is } k\} \Rightarrow B_k = A_k \setminus A_{k-1} \\
    \therefore& \abs{B_k} = \abs{A_k} - \abs{A_{k-1} } \\
    \Rightarrow& \Prob{B_k} = \frac{(k+1)^n - k^n}{10^n}
\end{align*}

\subsection{Coincident Birthdays}
Suppose $n$ people in a room. Find $\Prob(n) = $ probability that two have the same birthday.
Assume birthdays are numbered $\{1, 2, \cdots, 365\}$ and all equally likely.
\begin{align*}
    \Omega = \{1, 2, \cdots, 365\}^n \\
    A = \{\text{all different}\} \Rightarrow \abs{A} = \frac{365!}{(365 - n)!}
\end{align*}
So $\Prob(\text{two the same}) = 1 - \frac{365 \times 364 \times \cdots \times (365 -n +1)}{365^n}$
In fact,
\begin{align*}
    \Prob(22) = 0.476 \\
    \Prob(23) = 0.507
\end{align*}
So if $n \geq 23$ we are more than likely to have two with the same birthday.

\end{document}